\section{Sparse-deletion stability of the sampled quotient frame}
\label{sec:tpc51-mask-robustness}

We now leave the static mask intact and measure its effect directly on
the sampled Gram.  This is the order fixed in TPC--50; no
Fourier--Ferrers--gcd layer of a projective mask representation is
declared to be a physical Mellin cell.

Let \(\cR_X\) be the complete prime--squarefree row universe and let
\(M_X(\alpha,\gamma)\in\{0,1\}\) be the literal mask.  We use only the
inherited degree defect
\begin{equation}
 \sup_\alpha\#\{\gamma\in\cR_X:M_X(\alpha,\gamma)=0\}
 \leq\epsilon_X|\cR_X|,
 \label{eq:tpc51-mask-degree-hypothesis}
\end{equation}
and the same bound after interchanging the two row systems.  In the TPC
application,
\begin{equation}
 \epsilon_X\ll_H
 L^{-1+o(1)}+D^{-1+o(1)}+X^{-\kappa_0+o(1)}
 \label{eq:tpc51-inherited-mask-defect}
\end{equation}
by the complete-universe degree theorem inherited through TPC--49/50
\citep{WangTPC49,WangTPC50}.
For this TPC application we retain the hypotheses under which that
degree theorem was proved: \(L>4D\), and the row product
\(m=\ell d\) determines the prime--squarefree row.  The abstract masked
frame below applies whenever \eqref{eq:tpc51-mask-degree-hypothesis}
holds; the citation does not furnish that hypothesis outside this
complete-universe regime.

The elementary Nikolskii bound needed below is
\begin{equation}
 \sup_{(x,y)\in[0,T]^2}\|P_a(x,y)\|_{\cH}^2
 \leq R_K\sum_{k\in\Lambda_K}\|a_k\|_{\cH}^2,
 \label{eq:tpc51-nikolskii}
\end{equation}
which follows from Cauchy--Schwarz and remains valid for coefficients in
any Hilbert space \(\cH\).

\begin{theorem}[One-sided masked log-frame]
\label{thm:tpc51-masked-frame}
Fix the admissible row cell, base weights, and degree
\(K\leq(\log X)^A\) of \cref{thm:tpc51-polylog-fourier-frame}.  For
each source row \(\alpha\), let
\begin{equation}
 P_\alpha(x,y)=\sum_{k\in\Lambda_K}
 c_{\alpha,k}
 \exp\!\left(\frac{2\pi i}{T}(k_1x+k_2y)\right),
 \qquad c_{\alpha,k}\in\cH,
 \label{eq:tpc51-row-dependent-polynomial}
\end{equation}
with only finitely many nonzero source rows.  Then
\begin{align}
 &\sum_\alpha\sum_{\gamma\in\cR_X(\cA)}
 M_X(\alpha,\gamma)w_\gamma
 \|P_\alpha(x_\gamma,y_\gamma)\|_{\cH}^2
 \notag\\
 &\qquad\geq E_X
 \left(c_A-C_{A,\cA}\epsilon_XR_K\right)
 \sum_\alpha\sum_{k\in\Lambda_K}\|c_{\alpha,k}\|_{\cH}^2,
 \label{eq:tpc51-masked-lower-frame}
\end{align}
where \((x_\gamma,y_\gamma)\) are the logarithmic coordinates from
\eqref{eq:tpc51-empirical-measure}.  The corresponding upper bound is
\begin{equation}
 \sum_\alpha\sum_{\gamma\in\cR_X(\cA)}
 M_X(\alpha,\gamma)w_\gamma
 \|P_\alpha(x_\gamma,y_\gamma)\|_{\cH}^2
 \leq C_AE_X\sum_{\alpha,k}\|c_{\alpha,k}\|_{\cH}^2.
 \label{eq:tpc51-masked-upper-frame}
\end{equation}
In particular the lower bound remains positive whenever
\(\epsilon_XR_K=o(1)\).  The assertion is uniform in arbitrary
source-row concentration because the estimate is applied before the
sum over \(\alpha\).
\end{theorem}

\begin{proof}
The fixed admissible cell has positive density in the complete row
universe and its base weights satisfy
\begin{equation}
 E_X\asymp_{q,H,\cA}|\cR_X|(\log L)^{2e},
 \qquad
 \max_{\gamma\in\cR_X(\cA)}w_\gamma
 \ll(\log L)^{2e}.
 \label{eq:tpc51-base-flatness}
\end{equation}
For a fixed \(\alpha\), the contribution removed from the unmasked
quadratic form is therefore at most
\begin{align}
 &\#\{\gamma:M_X(\alpha,\gamma)=0\}
 \max_\gamma w_\gamma\,
 \|P_\alpha\|_\infty^2
 \notag\\
 &\qquad\ll_{q,H,\cA}
 \epsilon_XE_XR_K
 \sum_{k\in\Lambda_K}\|c_{\alpha,k}\|_{\cH}^2,
 \label{eq:tpc51-deleted-bank-energy}
\end{align}
where \eqref{eq:tpc51-mask-degree-hypothesis},
\eqref{eq:tpc51-base-flatness}, and \eqref{eq:tpc51-nikolskii} were
used in order.  Subtract this from the unmasked lower bound in
\eqref{eq:tpc51-unmasked-fourier-frame}, then sum in \(\alpha\).
Deleting nonnegative terms can only improve the unmasked upper bound,
which proves \eqref{eq:tpc51-masked-upper-frame}.
\end{proof}

\begin{corollary}[Fixed-shape quotient under the literal mask]
\label{cor:tpc51-masked-fixed-shapes}
Let \(W_1,\ldots,W_R\in C([0,T]^2)\) be a fixed family whose nonzero
quotient
classes are linearly independent in
\(L^2(e^{x+y}dxdy)\).  For every source row let
\(F_\alpha=\sum_{r=1}^R c_{\alpha,r}W_r\), with coefficients in an
arbitrary Hilbert space \(\cH\), and assume either finite source support
or \(\sum_{\alpha,r}\|c_{\alpha,r}\|_{\cH}^2<\infty\).  After
discarding exact zero classes
and choosing an independent quotient basis, there are constants
\(c,C>0\) such that
\begin{align}
 cE_X\sum_{\alpha,r}\|c_{\alpha,r}\|_{\cH}^2
 &\leq
 \sum_{\alpha}\sum_{\gamma\in\cR_X(\cA)}
 M_X(\alpha,\gamma)w_\gamma
 \|F_\alpha(x_\gamma,y_\gamma)\|_{\cH}^2
 \notag\\
 &\leq CE_X\sum_{\alpha,r}\|c_{\alpha,r}\|_{\cH}^2
 \label{eq:tpc51-fixed-shape-masked-frame}
\end{align}
for all sufficiently large \(X\).
\end{corollary}

\begin{proof}
By \cref{cor:tpc51-fixed-physical-family}, the unmasked Gram has a
uniform positive least eigenvalue.  Because the family is fixed,
\[
 \sup_{[0,T]^2}\left\|\sum_r c_rW_r\right\|_{\cH}^2
 \leq C_W\sum_r\|c_r\|_{\cH}^2.
\]
Repeating \eqref{eq:tpc51-deleted-bank-energy} gives an
\(O_W(\epsilon_X)\) perturbation, which is absorbed for large \(X\).
\end{proof}

\subsection{Canonical Mellin interpretation and the missing bridge}

The real-axis row powers in the provenance formula become, after
removing global unit phases,
\begin{equation}
 \ell^{it_1}d^{it_2}
 =L^{it_1}D^{it_2}
 e^{i(t_1x_\ell+t_2y_d)}.
 \label{eq:tpc51-log-phase-identification}
\end{equation}
For one declared fixed scalar function of \((x_\ell,y_d)\), integrating
these phases against its Mellin transform reconstructs that function
exactly; see \cref{thm:tpc51-canonical-mellin-direction}.  Such a fixed
finite scalar family is governed by
\cref{cor:tpc51-masked-fixed-shapes}.  This statement is not yet an
identification of the actual orbit-dependent TPC family with that
scalar model.

Indeed, the provenance factors include expressions such as
\(W(m_\gamma j/X)\) and \(\Psi(d_\gamma j/K)\).  With
\(z_j=j/J\), their row-log components have the schematic form
\begin{equation}
 W(c_Xe^{x+y}z_j),
 \qquad
 \Psi(c'_Xe^yz_j),
 \label{eq:tpc51-orbit-translate-shapes}
\end{equation}
for bounded scale factors on a fixed packet.  Retaining the orbit label
therefore produces an \(X\)-dependent Hilbert-valued row function, not
automatically a fixed list of scalar functions on \([0,T]^2\).

\begin{proposition}[Distinct log translates are independent]
\label{prop:tpc51-translate-independence}
Let \(0\ne w\in C_c^\infty(\R)\), and let
\(a_1,\ldots,a_M\in\R\) be distinct.  Then the functions
\(w(\,\cdot+a_1),\ldots,w(\,\cdot+a_M)\) are linearly independent on
\(\R\).
\end{proposition}

\begin{proof}
If \(\sum_jc_jw(s+a_j)=0\), Fourier transformation gives
\[
 \widehat w(\xi)\sum_jc_je^{ia_j\xi}=0.
\]
The transform \(\widehat w\) is not identically zero and hence is
nonzero on some open interval.  The exponential polynomial on the
right vanishes on that interval, so it vanishes identically.  Taking
its first \(M\) derivatives at one point gives a Vandermonde system in
the distinct \(a_j\), and therefore every \(c_j=0\).
\end{proof}

Applied to \(w(s)=W(e^s)\), this proposition shows that smoothness alone
cannot collapse the unrestricted orbit translates in
\eqref{eq:tpc51-orbit-translate-shapes} to a fixed-rank scalar family.
Restriction to the physical log box, support losses, and aggregation in
the orbit Hilbert space require a separate uniform argument.

We record exactly what such an argument must supply.  For a family
\(H=(H_\alpha)_\alpha\) of \(\cK_X\)-valued functions on the log box,
write
\begin{align}
 \mathcal Q_X^0(H)
 &:=\sum_\alpha\sum_{\gamma\in\cR_X(\cA)}
 w_\gamma\|H_\alpha(x_\gamma,y_\gamma)\|_{\cK_X}^2,
 \notag\\
 \mathcal Q_X^M(H)
 &:=\sum_\alpha\sum_{\gamma\in\cR_X(\cA)}
 M_X(\alpha,\gamma)w_\gamma
 \|H_\alpha(x_\gamma,y_\gamma)\|_{\cK_X}^2.
 \label{eq:tpc51-bridge-quadratic-forms}
\end{align}

\begin{proposition}[Actual-provenance bridge criterion]
\label{prop:tpc51-actual-bridge-criterion}
Suppose the coefficientwise actual row field, with source and residual
labels retained, has an exact representation
\[
 H_{\alpha,X}^{\rm act}(x,y)
 =\sum_{r=1}^{R_X}c_{\alpha,r,X}F_{r,X}(x,y),
 \qquad F_{r,X}\in C([0,T]^2;\cK_X),
\]
after quotienting its exact relations.  Let \(\mathcal C_X(c)\) be the
declared normalized coefficient energy, including the factor \(E_X\)
and the fixed shape normalizations.  Assume one of the following two
uniform alternatives.
\begin{enumerate}[label=\textup{(G\arabic*)},leftmargin=3.2em]
 \item The actual vector-valued family has a uniform normalized Gram and
 sparse-deletion estimate:
 \[
  c_0\mathcal C_X(c)\leq\mathcal Q_X^0(H(c))
  \leq C_0\mathcal C_X(c),
  \qquad
  \mathcal Q_X^0(H(c))-\mathcal Q_X^M(H(c))
  =o(\mathcal C_X(c))
 \]
 uniformly on the declared quotient subspace, for fixed
 \(0<c_0\leq C_0<\infty\).
 \item There is a retained log-Fourier bank \(P_X(c)\) of degree
 \(K_X\leq(\log X)^A\), with bank coefficient energy comparable to
 \(\mathcal C_X(c)\), such that the tail is uniformly relative:
 \[
  \mathcal Q_X^0(H(c)-P_X(c))=o(\mathcal C_X(c)).
 \]
\end{enumerate}
Finally suppose the actual weighted continuous Mellin envelope satisfies
the norm comparison
\begin{equation}
 \bigl(\cE_{1,B}^L(\one)\bigr)^2
 \leq X^{o(1)}\mathcal C_X(c^{\rm act}).
 \label{eq:tpc51-actual-envelope-bridge}
\end{equation}
for a nonzero actual quotient direction with
$\mathcal C_X(c^{\rm act})>0$.  Then, for all sufficiently large $X$, the
masked denominator below is positive and the coefficientwise,
pre-terminal actual direction has zero positive polynomial
representation loss:
\begin{equation}
 \frac{\bigl(\cE_{1,B}^L(\one)\bigr)^2}
 {\mathcal Q_X^M(H_X^{\rm act})}
 \leq X^{o(1)}.
 \label{eq:tpc51-conditional-actual-bridge}
\end{equation}
This conclusion retains the source and residual labels and makes no
statement about the later coherent prime or grouping maps.
\end{proposition}

\begin{proof}
Under \textup{(G1)}, subtraction gives
\(\mathcal Q_X^M(H(c))\geq(c_0-o(1))\mathcal C_X(c)\).
Under \textup{(G2)}, \cref{thm:tpc51-masked-frame} gives the same lower
bound for \(P_X(c)\), while
\(\mathcal Q_X^M(H-P_X)\leq\mathcal Q_X^0(H-P_X)=o(\mathcal C_X(c))\).
The triangle inequality in the masked \(L^2\) space transfers the lower
bound from \(P_X\) to \(H\).  Combining either lower bound with
\eqref{eq:tpc51-actual-envelope-bridge} proves
\eqref{eq:tpc51-conditional-actual-bridge}.
\end{proof}

Neither alternative nor \eqref{eq:tpc51-actual-envelope-bridge} is
verified here for the actual orbit-translate family
\eqref{eq:tpc51-orbit-translate-shapes}.  The fixed scalar corollary and
the retained finite-bank theorem prove model cases of the criterion;
they do not prove the actual bridge by relabelling.

Most importantly, the source label \(\alpha\) remains outside the
opposite-row norm in \eqref{eq:tpc51-masked-lower-frame}.  If a later
operation coherently sums different source rows before taking a norm,
new cross-row Gram terms appear.  The theorem neither discards nor
estimates those terms.
